\minitopic{专题研讨：概率模型、校准与决定}形式认识论把“较相信”“证据更强”转化为可比较结构，却不把认识论缩成算术。任何概率都相对于命题划分、背景信息和模型假设。算式正确而事件定义错误，可能比没有数字更具误导性。本节的目标是把计算置于完整工作流：先规定问题，再更新信念，检查校准，最后引入价值作出决定。

\minitopic{后验概率与似然不能交换}检测阳性时，人们容易把“患者患病的概率”误当成“患者患病时检测阳性的概率”。前者是 \(P(D\mid +)\)，后者是 \(P(+\mid D)\)。贝叶斯公式之所以重要，不只是提供一个换算，而是迫使我们同时保留基础率 \(P(D)\)、真阳性率和假阳性率。罕见疾病即使检测敏感，也可能产生大量假阳性。 似然比比较证据在两个假设下的预期程度。若 \(e\) 在 \(h\) 下远比在非 \(h\) 下可能，\(e\) 支持 \(h\)；但“强支持”不等于“高后验”。一个极低先验的假设可以被证据大幅提高，结果仍不够高。报道应同时说明变化幅度和最终水平，避免用“概率增加三倍”掩盖从万分之一到万分之三。 自然频率把抽象百分比转换成计数。在一万人中先画患病与未患病人数，再填检测结果，许多方向错误会显现。频率图还揭示分母选择：是所有检测阳性者、所有患病者，还是所有被送检者？形式训练的首要德性是不给分母消失的机会。

\minitopic{先验从何而来}先验并非纯粹主观任性。它可以来自基准频率、先前研究、对称性或信息受限下的约束；也可能表达合理分歧。争论先验时应问：参考类如何选择，旧研究与当前对象是否可比，先验对结论有多敏感？如果一组合理先验都导向相近结论，结果稳健；若轻微变化就逆转，报告应突出先验敏感性。 所谓“无信息先验”也依赖参数化。在发生概率上均匀，不等于在对数几率上均匀。选择不可能完全脱离建模目的，因此透明说明比宣称中立更重要。客观约束可以缩小合理范围，却未必给出唯一数字。\parencite{howson2006} 先验还可能携带历史偏差。犯罪风险模型用过去逮捕率作基础率，逮捕率又受巡逻分布影响；把它当作自然频率会把制度选择写进预测。形式上合法的更新不保证输入具有正当解释。模型说明应区分目标事件、可观测代理与生成数据的社会过程。

\minitopic{条件独立与重复证据}连续获得两项证据时，只有在相关条件下独立，才能简单相乘似然。两次检测使用同一污染样本、两位专家读到同一初步报告、两篇论文使用同一数据库，都可能共享错误。忽略相关性会重复计算证据，使后验虚高。 判断相关性应针对假设条件。两个症状在人群中相关，不表示在已知疾病状态后仍相关；两名证人平时互不认识，若事后共同观看新闻，也不再独立。贝叶斯网络用有向边表示依赖，帮助追问哪些共同原因被遗漏。图并不自动保证因果正确，但能把隐藏假设暴露给批评。 证据也可能相互削弱或增强。两项各自模糊的测量若误差方向相反，联合后更精确；两个几乎相同的来源只增加很少信息。报告“共有五项证据”没有意义，除非说明它们的联合结构。数量是证据的属性之一，不是支持度本身。

\minitopic{校准、准确性与评分}一个预报者若对所有标为 \(70\%\) 的事件中约七成最终发生，称为在该组上校准。校准评价长期匹配，不保证个案正确，也不保证预测足够精细。总把晴雨各报 \(50\%\) 在某些数据上可以校准，却没有区分能力。评价应结合校准、锐度、分辨率和适用范围。 恰当评分规则让诚实报告真实置信度在期望上最有利。它们可以比较概率预测而不只比较二元命中。\parencite{joyce1998} 但分数仍依赖事件定义与样本选择。平台若只展示已经解决的预测题，或在临界时更改结算标准，漂亮分数会误导。形式指标需要可审计的数据治理。 校准还应分群体检查。总体表现良好可能掩盖在某年龄、地区或设备上的系统偏差。细分过度又会因样本太小而不稳定。合理做法是预先选择与故障机制相关的亚组，给出不确定区间，并在新数据上复验，而不是事后寻找最有利切片。

\minitopic{概率并不直接给出行动}行动需要把概率与结果价值结合。期望效用以每种结果的概率乘其效用，再比较方案。若误治副作用很小，较低患病概率也可支持治疗；若手术不可逆且危险，可能需要更高阈值。阈值不是“知道”的简单数值定义，而是具体决定中的损益折中。\parencite{jeffrey1983} 效用也不是金钱的同义词。生命、隐私、公平、可逆性和程序权利可能进入评价。有些限制不适合完全换算成单一数值，例如不得未经同意把个体置于高风险实验。决策分析可以展示权衡，却不能替代伦理论证。 信息本身也有价值。一次追加检测的价值取决于它改变行动的概率、检测成本和延误风险。若无论结果如何都采取同一方案，检测可能没有决策价值；但它仍可能具有解释、安慰或研究价值。区分认识价值与决策价值可防止“不能改变行动，所以不值得知道”的草率结论。

\minitopic{模型比较与证据欠定}同一数据可能由多个模型解释。更复杂模型通常拟合更好，却可能过拟合；更简单模型易推广，却可能漏掉真实机制。模型比较需在拟合、复杂度、预测性能和解释目标间权衡。交叉验证测试未见数据表现，不能证明部署环境永远相同；信息准则压缩复杂度代价，也不能代替领域知识。 证据欠定不意味着任意选择。模型可能在旧数据上近似等价，却对新干预、边界条件或辅助变量给出不同预测。研究者可设计区分性实验，寻找双方预期差异最大的观察。若当前无法区分，报告模型不确定性比挑一个最顺眼的故事更诚实。 形式模型的简化必须可追踪。把连续风险二分为高低、把未知值当作阴性、假设损失对称，都会改变结论。每个简化应附有适用理由与敏感性检查。模型不是世界的缩小复制，而是为特定问题保留某些关系的工具。

\minitopic{比较视野：度量与审势}古典实践推理常强调审时、度势与权衡轻重，这与概率决策在功能上相交：行动者不能只问命题真伪，还要判断情势与后果。但这种相交不表示古典“权”就是现代期望效用。前者嵌在德性、礼义和具体判断中，未必接受一切价值可通约；后者提供明确的一致性结构。 比较可以形成双向批评。形式方法要求古典式权衡说明哪些因素改变结论，避免“因时制宜”成为无约束直觉；实践智慧传统则提醒形式模型，事件划分和效用赋值本身需要情境判断。目标不是让一方翻译掉另一方，而是检查各自遗漏：一种可能缺少可计算约束，另一种可能遮蔽价值来源。

\minitopic{综合案例：两阶段癌症筛查}某疾病在人群中的基础率为 \(1\%\)。初筛敏感度 \(90\%\)、特异度 \(95\%\)；复筛使用同一生物标志物，敏感度 \(95\%\)、特异度 \(90\%\)。机构把两次阳性当作独立证据，宣称患病概率超过 \(75\%\)。但两项检测受同一种炎症影响，条件独立不成立。 分析应先用自然频率计算独立假设下的结果，再给相关性设定一组合理范围，展示后验敏感性。随后区分三项决定：是否告知患者、是否做侵入式活检、是否立即治疗。三者损益不同，不能共用一个概率阈值。报告还要说明基础率是否适用于有症状的送检人群。 若医院拥有过往联合检测数据，可直接估计两次结果的联合条件概率；若样本不足，应给区间而非伪精确数。最后检查校准：过去的“双阳性”患者中实际患病比例是多少，各亚组是否一致？一个成熟答案不会止于单一后验，而会把模型不确定性转化成进一步取证与决策方案。

\begin{tcolorbox}[breakable,colback=white,colframe=blue!30!black,title={专题研读任务},fonttitle=\bfseries]
找一则带百分比的公共报道，重建事件、分母、先验、似然与决策阈值。至少提出一个相关性问题和一个参考类问题，并把原报道改写成同时呈现数值与假设的版本。
\end{tcolorbox}

\index{似然比}
\index{先验概率}
\index{条件独立}
\index{校准}
\index{恰当评分规则}
\index{期望效用}
